% ===========================================================================
\chapter{The Training Objective Decides Whether Attention Weights Stay Symmetric}
\label{chap:attention}
\margintoc
% ===========================================================================

\begin{publicationbasis}{Published}
Saponati, M., Sager, P.\,J., Vilimelis Aceituno, P., Stadelmann, T., \&
Grewe, B.\,F. (2025). The underlying structures of self-attention: Symmetry,
directionality, and emergent dynamics in Transformer training.
\textit{Proceedings of the 42nd International Conference on Machine Learning
(ICML)}. \cite{saponati_underlying_2025}
\end{publicationbasis}

%% Taken from the authorship contributions section of the arXiv version of
%% this paper, imported_papers/self_attention_geometry_project_ICML/arXiv/
%% main-arxiv.tex, which the ICML camera ready drops. It is the authors' own
%% statement and not a reconstruction.
\contributionstatement

Matteo Saponati and the thesis author are shared first authors, and Matteo
Saponati led the work. The paper's own statement of contributions reads:
conceptualisation by Matteo Saponati, Thilo Stadelmann and Benjamin F. Grewe;
mathematical analysis by Matteo Saponati and Pau Vilimelis Aceituno; analysis
of the pretrained models by Matteo Saponati and the thesis author; experiments
by the thesis author; writing by Matteo Saponati, the thesis author and
Benjamin F. Grewe; supervision by Thilo Stadelmann and Benjamin F. Grewe.

\section{Chapter Introduction}

\Cref{chap:grokking} varied the learning rule with everything else fixed and
found that the rule sets how much of a network carries the computation. The
natural next question is whether that is a peculiarity of the one component
that was varied, or whether other components of the training setup also leave
structure of their own in a trained network. This chapter answers that for a
second component, in a different substrate and at a scale the previous setting
cannot reach.

The design is the mirror image of the previous chapter's. Backpropagation is
held fixed and the training objective is varied, between predicting a masked
token from the whole sequence and predicting each token from its predecessors.
The object measured is the product of the query and key matrices of an
attention head, which is the bilinear form through which the two matrices act
on the embeddings. The structure in question is whether that form stays close
to symmetric, or whether it becomes directional, meaning that its columns of
high norm come to outnumber its rows of high norm.

What this chapter contributes to the thesis argument is one claim and not
another. It shows that a second component of the training setup leaves its own
structure in the weights, which is what makes the previous chapter's finding a
case of something rather than a curiosity. It is not a second measurement of
the circuit. No participation ratio over neurons is computed here, the object
is a weight matrix rather than a layer, and the property measured is an
asymmetry between the rows and the columns of that matrix rather than a
concentration of the computation on a few of either. The variable is the
objective, and in a decoder part of that variation is carried by the mask over
which positions may attend to which, which belongs to the architecture.

The paper's own term for this invites a misreading and needs pinning down. It
speaks of \emph{symmetry} and \emph{directionality} of the attention
weights, which are properties of the query key product under transposition,
and of \emph{column dominance} when columns of high norm are more common than
rows of high norm. Column dominance is columns against rows. It is not a few
columns dominating the others, and it is not the thesis's circuit measure,
which counts participating units in a layer. The chapter's $W_{qk}$ is
the product written $W_Q^\top W_K$ elsewhere in the literature.

% ---------------------------------------------------------------------------
% The published paper, included verbatim. Do not edit imported_papers/.
% ---------------------------------------------------------------------------

% The paper keeps its section headings and their labels in its own root
% file, which the chapter replaces, so both are restored here. Written
% inside the block, the labels pick up the paper's slug prefix, which is
% what the paper's own cross references expect.
\begin{importedpaper}{attention}
\section{Autoregressive and bidirectional training leads to directional and symmetric weight updates}
\label{sec-results-math}
\importsection{ICML2025/sections/results-math}

\section{Symmetric and directional structures are predominant in Transformer models}
\label{sec-results-pretrained-models}
\importsection{ICML2025/figures/figure-pretrained-models-language}
\importsection{ICML2025/sections/results-pretrained-models}

\section{Experiments}
\label{sec-results-experiments}
\importsection{ICML2025/sections/results-experiments}
\end{importedpaper}

\section{Chapter Summary}

Every gradient step on a Transformer's attention weights accumulates over
whichever ordered token pairs the objective allows, and the two standard
objectives allow different ones. Bidirectional training enters every pair in
both orders, once with each token as the masked target, so the contributions
cancel antisymmetrically and the query key product stays close to symmetric.
Autoregressive training makes each token the target exactly once and the
context of every token after it, so the sum is one sided and the product
becomes directional, with columns of high norm outnumbering rows of high
norm. The derivation is not assumption free. It drops
the softmax, omits layer normalisation and biases, takes one small step of
plain gradient descent where the experiments use an adaptive optimizer, and
the directionality result needs the token embeddings to be anisotropic and a
corpus statistic to favour context over target.

The prediction is borne out in trained models of both kinds. Twelve layer
encoders and decoders trained under matched conditions on three corpora
separate cleanly on the two measures across training, and the same separation
appears in released models, in ModernBERT on one side and in GPT, LLaMA3 and
Mistral on the other, and across text, vision and audio. The matched runs are
the controlled evidence and the released models show that the effect survives
at the scale the field works at.

One consequence of the derivation is practical. If bidirectional training
drives the query key product towards symmetry, then initialising it symmetric
starts the network closer to where training is going, and it does. Encoder
models initialised that way reach the final loss of a randomly initialised
model as much as $75\%$ sooner at four layers and $35\%$ sooner at twelve, and
the spread across corpora is wide: at twelve layers the three runs give
$35\%$, $20\%$ and nothing at all. The paper runs no matching comparison on
decoders, so whether the same initialisation would help or hurt there is
untested.

Two components of the training setup now have a structure attached to them.
The learning rule settles how much of a layer carries a computation, the
objective settles whether a query key product stays symmetric, and the task
fixes neither.

Both chapters so far have measured a trained network and drawn a conclusion
from the measurement. \Cref{chap:bounded} asks what has to be true of such a
measurement before the conclusion follows.

